\section{Introduction}
\label{sec:intro}

Generative models have entered human-robot interaction (HRI) by two doors.
Through the first, they have become parts of robots. Large language models
(LLMs) write what social robots say
\cite{williams2024scarecrows,kim2024understanding,wang2026how}; they and
their multimodal successors, vision-language models (VLMs), interpret what
people say and do, choose what the robot should do next, and generate
gestures, facial expressions, and programs
\cite{mahadevan2024generative,wang2024lami}. Vision-language-action models
(VLAs) go further and map camera images and spoken instructions directly to
motor commands \cite{zitkovich2023rt2,kim2025openvla,black2025pi0}. Through
the second door, less noticed, the same models have become parts of the
research apparatus: they code transcripts, score video, summarise interviews,
judge the quality of other models' outputs, and simulate the users on whom
protocols are piloted \cite{gilardi2023chatgpt,ziems2024can,zheng2023judging}.

The community has no shared standard for evaluating either use. The question
was put plainly in a recent meeting of HRI researchers concerned with user
studies and with human-robot teams. The wider AI community offers practices,
such as using a model as a judge of another model's outputs, or validating
outputs against human judgement, but HRI needs standards of its own, because
what counts as adequate evaluation depends on how the model is used.
Comparing a model's natural-language outputs with those a human would have
produced is a way of establishing that the outputs are relevant, safe,
transparent, and correct. Evaluating a model whose embeddings feed a
downstream component looks instead like evaluating a classifier. It is hard,
the meeting concluded, to see a one-size-fits-all solution.

This article takes that conclusion as its starting point and turns it into a
proposal. We argue that a single standard for ``LLMs in HRI'' is neither
achievable nor desirable, and that the difficulty dissolves once the unit of
evaluation is changed. What should be evaluated is not the model but the
\emph{role} the model plays in the robot's architecture or in the research
pipeline. Roles are few and stable even as models change every few months,
and each role already has an evaluation tradition with decades of
methodological development behind it. A model that speaks to a person
inherits the tradition of human evaluation of generated text
\cite{vanderlee2021human,howcroft2020twenty}. A model that perceives a
person's state inherits the traditions of classifier evaluation and of
behavioural coding with multiple raters \cite{hallgren2012computing}. A model
that moves an arm near a person inherits the tradition of trial-based policy
evaluation \cite{kressgazit2024robot}. A model that codes the data of a study
inherits the tradition of content analysis, in which a coder, human or
otherwise, is trusted only to the extent that their agreement with other
coders has been demonstrated \cite{krippendorff2019content}. The work of
standard-setting is then largely one of translation: identifying the role,
finding the tradition, and adapting its requirements to the conditions of HRI.

Those conditions are what make HRI's problem distinct from the problem that
natural language processing (NLP) is already working on. First, a robot's
outputs are embodied and situated. An utterance is delivered in a voice, by a
body, in a place, to a particular person, within a conversational pause; a
response that would be rated as excellent on a screen can fail as speech
because it is too long, arrives two seconds late, or refers to something the
robot cannot see \cite{skantze2021turn,irfan2025between}. Second, outputs can
become physical actions, so that an error is not only a wrong sentence but a
movement near a body \cite{robey2025jailbreaking,hundt2025llm}. Third, HRI
studies are small. Samples of a few dozen participants and a few hundred
utterances are typical, which rules out the large test sets on which
meta-evaluation in NLP relies and makes the precision of every estimate a
live concern \cite{leichtmann2022crisis}. Fourth, many HRI applications
involve children, older adults, and people seeking support, populations for
whom an inadequate output has consequences that a benchmark score does not
express \cite{abbasi2025robotled,irfan2024recommendations}. Fifth, the field's
primary instrument is the user study, and it is tempting to let the user
study stand in for evaluation of the model. We will argue that it cannot.

There is also evidence that current practice falls short of what the
inherited traditions would accept. A systematic review of 86 LLM-driven HRI
studies found that questionnaires were used in 70 of them and metrics
specific to the model's performance in 30, and characterised the field as
largely exploratory, with widely varying setups, methods, and metrics
\cite{wang2026how}. Reporting is uneven: some studies give a dated model
snapshot, sampling parameters, and full prompts, and others name only a
product, which is why reporting guidelines for LLMs in HRI have recently
been proposed \cite{williams2024scarecrows,matuszek2026reporting,
ashok2026why}. HRI has in
the past identified and repaired weaknesses of this kind, in the reporting of
Wizard-of-Oz studies \cite{riek2012wizard}, in the statistical treatment of
rating scales \cite{schrum2023concerning}, and in replicability more broadly
\cite{gunes2022reproducibility,leichtmann2022crisis}. Generative models are
the next such case, and they compound the earlier ones. LLMs in robot
architectures have been aptly described as scarecrows, black-box modules
that stand where the wizard used to stand \cite{williams2024scarecrows},
and the scarecrow's brain is usually rented from a provider that may change
it without notice \cite{chen2024how}.

\paragraph{Contributions.} The article makes four contributions.
\begin{enumerate}[leftmargin=1.6em]
\item A \emph{role-sensitive framework} (Section~\ref{sec:framework}) that
distinguishes nine roles of generative models in HRI research, four levels at
which evidence is gathered, five kinds of claim, and three tiers that scale
the expected evidence to exposure and consequence.
\item Twenty-four \emph{best practices}: ten that apply to any study
(Section~\ref{sec:crosscutting}) and fourteen that are specific to a role
(Section~\ref{sec:rolepractices}), each traced to the tradition in NLP,
psychology, or robotics from which it is adapted.
\item A description of \emph{emerging practices} (Section~\ref{sec:emerging})
that are new or newly feasible, several of them first tried in HRI, with an
assessment of what each can and cannot establish.
\item Practical instruments: five \emph{worked examples}
(Section~\ref{sec:examples}), \emph{precision-planning tables} for validation
subsets and physical trials (Appendix~\ref{app:precision}), and a
\emph{reporting checklist} with questions for reviewers
(Section~\ref{sec:reporting}).
\end{enumerate}

\paragraph{Scope and approach.} This is an integrative methodological review
with recommendations, not a systematic review. We did not attempt to
enumerate and code every HRI study that uses a generative model; recent
systematic reviews do that \cite{wang2026how,nikolova2026systematic}. We read
across three literatures (evaluation methodology in NLP and machine learning;
reliability, validity, and research practice in psychology and content
analysis; policy evaluation and safety in robotics) for practices that answer
problems HRI studies now face, and we examined recent HRI work, including our
own, for practices that are emerging locally. The selection reflects the
author's position in social and affective robotics, where the Speaker and
Perceiver roles and vulnerable populations are most prominent; we have tried
to give physical interaction and action models equal weight, and we return to
this limitation in Section~\ref{sec:discussion}. Throughout, ``model'' means a
pretrained generative or foundation model (LLM, VLM, or VLA) together with
its prompt and configuration, and ``evaluation'' means the gathering of
evidence about whether the model is adequate to its role, as distinct from
the evaluation of the robot or of the interaction as a whole. The
recommendations are proposals for discussion. A standard becomes one only
when a community adopts it, and Section~\ref{sec:discussion} suggests how that
process might proceed.
